\documentclass[10pt,a4paper]{amsart}
\usepackage[margin=19mm]{geometry}
\usepackage{amsmath,amssymb,mathtools}
\usepackage[scr=rsfs,cal=euler]{mathalfa}
\usepackage{xfrac}
\usepackage[unicode,hidelinks,bookmarks=false]{hyperref}
\newcommand{\ie}{i.e.\ }
\newcommand{\cf}{cf.\ }
\newcommand{\resp}{respectively\ }
\newcommand{\Subsection}{\subsection}
\providecommand{\nobreakdash}{\nobreak\hbox{-}}
\setlength{\emergencystretch}{2em}
\multlinegap0pt
\usepackage{bm}
\usepackage{bbm}
\usepackage{dsfont}
\usepackage{xspace}
\usepackage{cancel}
\usepackage{mathtools}
\usepackage{amsthm}
\makeatletter
\@ifundefined{theorem}{\newtheorem{theorem}{Theorem}}{}
\@ifundefined{proposition}{\newtheorem{proposition}[theorem]{Proposition}}{}
\makeatother
\newcommand{\fqm}{\mathbb{F}_{q^m}}
\newcommand{\fq}{{\mathbb F}_{q}}
\title[Delsarte duality on subspaces and applications to rank-metri]{Delsarte duality on subspaces and applications to rank-metric codes and q-matroids}
\author{Martino Borello, Olga Polverino, Ferdinando Zullo}
\date{Source: arXiv:2509.24409v1 (CC BY 4.0)}
\begin{document}
\maketitle

\noindent\emph{Source and licence.} Delsarte duality on subspaces and applications to rank-metric codes and q-matroids. Version arXiv:2509.24409v1 (CC BY 4.0), distributed under the Creative Commons Attribution 4.0 International licence, as recorded in the arXiv licence metadata for this exact version and shown on the abstract page https://arxiv.org/abs/2509.24409v1. Full rights notice: licenses/arxiv-2509.24409-CC-BY-4.0.txt.\par\smallskip

\noindent\textit{Editorial scope.} Counted unit: the Theorem whose source label is \texttt{None} in the article (source0.tex lines 830--833, source label \texttt{None}), delivered together with the article's own complete original proof of it (source0.tex lines 835--874, one \texttt{proof} environment of 40 lines). No mathematics is written, repaired or re-derived: statement and proof are the article's own text line for line. Required context retained uncounted: def:F\_complwe (lines 428--430); prop:decomposable (lines 441--450); thm: dual defect (lines 629--645); thm: =dual defect (lines 667--675). Every definition, display and label of those lines is retained unchanged. Omitted from this package and not redistributed: 1--427, 431--440, 451--628, 646--666, 676--829, 834--834, 875--1653. Nothing inside a retained excerpt is omitted or altered: every retained body is delivered exactly as the article's lines are written, so no omission receipt of any kind accompanies this package. Citations: the retained excerpts carry no citation command at all, so no bibliography entry of the article is transcribed and no reference is redistributed. Reference adaptation: none was needed and none was made; every reference inside the delivered unit resolves inside it, so no unresolved reference or citation is produced. Printed slips kept literal: no printed word, symbol, hypothesis or number of the counted statement or of its proof was altered. Layout and class: the article's own class is \texttt{amsart} with packages including amsmath, relsize, cite, color, xcolor, mathtools, ulem, hyperref, enumitem, amsthm, amssymb, geometry, caption, tikz-cd; the wrapper is the lane's pinned \texttt{amsart} editorial wrapper at 10pt on A4, which restates the theorem-like environments the retained text needs and adds the article's own macro definitions that the retained text needs (\texttt{\textbackslash fq}, \texttt{\textbackslash fqm}), with extra wrapper packages (bm, bbm, dsfont, xspace, cancel, mathtools). The delivered numbering is the unit's own. Assets: no retained span contains a figure environment, a graphics-inclusion command, a PDF-page inclusion, a table or a TikZ picture; no image file is copied into this package, the package declares no asset dependency and no image file is redistributed with it. Licence: version arXiv:2509.24409v1 of this article is distributed under the Creative Commons Attribution 4.0 International licence, as recorded in the arXiv licence metadata for this exact version and shown on the abstract page https://arxiv.org/abs/2509.24409v1. A full rights notice is delivered at licenses/arxiv-2509.24409-CC-BY-4.0.txt. Characters: every retained excerpt is pure ASCII, so the delivered engine, XeLaTeX with its default Latin Modern text font, reports no missing character.
  \begin{equation} \label{def:F_complwe}
      \mathbb{V}(k,q^m)=F_1\oplus \ldots \oplus F_t ,
  \end{equation}

\begin{proposition} \label{prop:decomposable}
 Let $U$ be in $\mathcal{L}_q(\mathbb{V}(k,q^m))$ with dimension $n$. Suppose that $U$ is decomposable of type $(\mathbf{k}, \mathbf{n})$, with $n=\sum_{i=1}^t n_i$, and let $F_i$ be, for $i\in[t]$, the components of the decomposition  (\ref{def:F_complwe}). Then the  following holds
 \begin{itemize}
    
     \item[1.] If $T$ is in $\mathcal{L}(\mathbb{V}(k,q^m))$ such that $T=\sum_{j=1}^{h}F_{i_j}\oplus T'$ for $h$ distinct components $F_{i_j}$ and $T'\leq_{\fqm} \sum_{l=1,\neq i_j}^{t}F_l$, then 
     \[w_U(T)=\sum_{j=1}^{h}w_U(F_{i_j})+w_U(T') \,\,\, \text{and} \,\,\, \varepsilon_U(T)=\sum_{j=1}^{h}\varepsilon_U(F_{i_j})+\varepsilon_U(T').\,\,\, \]
     \item[2.]  The minimum dimension of a subspace of $\mathbb{V}(k,q^m)$ with defect $n-k$ is $k$ (i.e. $t_s=k$) if and only if the components $F_i$ for every $i\in [t]$ are minimal with respect to its defect.
     \item[3.]  If $t_s=k$, then every subspace $T$ in $\mathcal{L}(\mathbb{V}(k,q^m))$ that is the sum of some components $F_i$  is minimal with respect to its defect.
 \end{itemize}
\end{proposition}

\begin{theorem} \label{thm: dual defect}
Let  $T$ be an $\fqm$-subspace of $\mathbb{V}(k,q^m)=\frac{\mathbb{V}(n,q^m)}{\Gamma}$ with $w_{U}(T)=h$ and such that $\langle T \cap U \rangle_{\fqm}=T$. Let $T=\frac{S_h^*+\Gamma}{\Gamma}$ where $S_h$ is an $h$-dimensional $\fq$-subspace of $W$. Let $T^d$ be the $\fqm$-subspace of $\mathbb{V}(n-k,q^m)=\frac{\mathbb{V}(n,q^m)}{\Gamma^\perp}$ defined as follows
\begin{equation} \label{eq:T^d}
T^d= \frac{(S_h^*)^\perp+\Gamma^\perp}{\Gamma^\perp}.   \end{equation} 
The following hold:
\begin{itemize}
    \item[\bf 1.] $\mathbb{V}(k,q^m)^d=\{ \underline 0\}$;
    \item[\bf 2.]  $\dim_{\fqm} (T^d)=n-k-\varepsilon_{U}(T)$;
    \item[\bf 3.] $w_{U^d}(T^d)\geq n-w_{U}(T)$;
    \item[\bf 4.] $\varepsilon_{U^d}(T^d)\geq k- \dim_{\fqm} (T)$;
    \item[\bf 5.] $(T^d)^d \subseteq T$ and $\varepsilon_{U}((T^d)^d)\geq \varepsilon_{U}(T)$;
    \item[\bf 6.] If $T'$ is an $\fqm$-subspace of $\mathbb{V}(k,q^m)$ such that
    \[T'\subseteq T \,\,\,\, \mbox{and} \,\,\,\, \langle T' \cap U \rangle_{\fqm}=T',\]
we have
\[T^d\subseteq T'^d.\]
\end{itemize}
\end{theorem}

\begin{theorem} \label{thm: =dual defect}
Let $T$ be an $\fqm$-subspace of $\mathbb{V}(k,q^m)=\frac{\mathbb{V}(n,q^m)}{\Gamma}$ with $w_{U}(T)=h$ such that $\langle T \cap U \rangle_{\fqm}=T$ and $\varepsilon_{U}(T)>0$. If $T$ is minimal with respect to its defect, then the inequalities in {\bf 3.}, {\bf 4.} and {\bf 5.} of Theorem \ref{thm: dual defect} are equalities, i.e.
\begin{itemize}  
    \item[\bf 1.] $w_{U^d}(T^d) =n-w_{U}(T)$;
    \item[\bf 2.] $\varepsilon_{U^d}(T^d)=k- \dim_{\fqm} (T)$;
    \item[\bf 3.] $(T^d)^d = T$.    
\end{itemize}
Moreover, $T^d$ is minimal with respect to its defect as well.
\end{theorem}

\begin{theorem}
    Let $\mathbf{k}=(k_1,\ldots,k_t)$,  $\mathbf{n}=(n_1,\ldots,n_t)\in\mathbb N^t$ with $k=k_1+\ldots+k_t$,  $n=n_1+\ldots+n_t$  and $k_i< n_i$ for $i\in [t]$. Consider a decomposable $\fq$-subspace $U$ of $\mathbb{V}(k,q^m)$ of dimension $n$  of type $(\mathbf{k}, \mathbf{n})$.
    We have that a Delsarte dual $U^d$ of $U$ is a decomposable subspace in $\mathbb{V}(n-k,q^m)$ of type $(\mathbf{n}-\mathbf{k}, \mathbf{n})$.
\end{theorem}

\begin{proof}
   Consider the representation of $U$ in the quotient model, for which we have $\mathbb{V}(k,q^m)=\frac{\mathbb{V}(n,q^m)}{\Gamma}$ and  $U=\frac{W+\Gamma}{\Gamma}$, where $W$ is a $n$-dimensional $\fq$-subgeometry of $\mathbb{V}(n,q^m)$. Since $U$ is decomposable of type $(\mathbf{k},\mathbf{n})$, we may write
   \[U=U_1\oplus \ldots \oplus U_t=\frac{S_1+\Gamma}{\Gamma}\oplus \cdots \oplus \frac{S_t+\Gamma}{\Gamma}\]
   and 
   \[\mathbb{V}(k,q^m)=F_1\oplus  \cdots \oplus F_t=\langle U_1 \rangle_{\fqm}\oplus \cdots \oplus \langle U_t\rangle_{\fqm}=\frac{S_1^*+\Gamma}{\Gamma}\oplus \cdots \oplus \frac{S_t^*+\Gamma}{\Gamma},\]
   where $S_i$ is an $n_i$-dimensional $\fq$-subspace of $W$.
   Also, since $W\cap \Gamma=\{\underline 0\}$, we have
   \[W=S_1\oplus \cdots \oplus S_t.\]
   Then it is easy to check that if $H_i=\sum_{j=1, j\neq i}^{t}S_j$, then 
   \begin{equation} \label{eq:W}
       W=\bar{S}_1\oplus \cdots \oplus \bar{S}_t
   \end{equation}
   where $\bar{S}_i=(H_i)^{\perp'}=\cap_{j=1, j\neq i}^{t} S_j^{\perp'}$ and $\dim_{\fq}(\bar{S}_i)=n_i$. 
   Also, if $I=\{i_1,i_2,\dots, i_s\} \subseteq [t]$, since $\sum_{j=1}^s S_{i_j}=\cap_{i\in [t]\setminus I} H_i $, we have  
   \begin{equation} \label{Eq:S partial sum}
      \cap_{j=1}^{s}S_{i_j} ^{\perp'}=\sum_{i\in [t]\setminus I}^{s}\bar{S}_i.
   \end{equation}
   Let $F'_i=\sum_{j=1, j\neq i}^{t}F_j$ for $i\in [t]$ and note that 
   \[F'_i=\sum_{j=1, j\neq i}^{t} \frac{S_j^*+\Gamma}{\Gamma}=\frac{(\sum_{j=1, j\neq i}^{t}S_j)^*+\Gamma}{\Gamma}=\frac{H_i^*+\Gamma}{\Gamma}.\]
   Clearly, the subspaces $F'_i$ have positive defect and, by {\it 3.} of Proposition \ref{prop:decomposable}, they are  minimal with respect to their defects. Hence, by Theorems \ref{thm: dual defect} and  \ref{thm: =dual defect} we have that 
   \[(F_i')^d=\frac{(H_i^{\perp'})^*+\Gamma^{\perp}}{\Gamma^{\perp}}=\frac{\bar{S}_i^*+\Gamma^{\perp}}{\Gamma^{\perp}},\]
   \[\dim_{\fqm}((F_i')^d)=n-k-\varepsilon_U(F_i')=n_i-k_i,\] 
   \[ w_{U^d}((F_i')^d)=n-w_U(F_i')=n_i,\] 
   \[\varepsilon_{U^d}((F_i')^d)=k-\dim_{\fqm}(F_i')=k_i,\] 
   and
   \[((F_i')^d)^d=F_i'.\]
   Also, $(F_i')^d$ is minimal with respect to its defect. By Equation (\ref{eq:W}) we get
   
   \begin{equation} \label{eq:decomposition_Ud}
       U^d=\frac{W+\Gamma^{\perp}}{\Gamma^{\perp}}=\frac{\bar{S}_1+\Gamma^{\perp}}{\Gamma^{\perp}}\oplus\cdots\oplus \frac{\bar{S}_t+\Gamma^{\perp}}{\Gamma^{\perp}}
   \end{equation}
   and hence, since $(U^d)^*=\mathbb{V}(n-k,q^m)$, we have
   \[\mathbb{V}(n-k,q^m)=(F_1')^d+\cdots +(F_t')^d.\]
   
and now comparing the dimensions we get
\begin{equation} \label{eq:decomposition_Vd}
\mathbb{V}(n-k,q^m)=(F_1')^d\oplus\cdots \oplus (F_t')^d,
\end{equation}
i.e. $U^d$ is decomposable of type $(\mathbf{n}-\mathbf{k}, \mathbf{n})$.
\end{proof}

\end{document}
